% TOP_PROBLEM: {"id": 30006757, "title": "Density of Hyperbolicity for Rational Maps", "category_id": 11, "category": {"id": 11, "name": "dynamical_systems", "display_name": "Dynamical Systems", "description": "Problems about long-term behavior of deterministic systems, Hamiltonian dynamics, and periodic orbits.", "slug": "dynamical-systems", "order_index": 11, "created_at": "2026-07-31T15:26:25.671Z"}, "status": "open"}
% FIRST_POSTED: 2026-09-27
% !TeX program = lualatex
% Compile twice: lualatex 30006757.tex
% All references and prose are embedded; no BibTeX database or external inputs.
\documentclass[11pt,a4paper]{article}
\usepackage[margin=24mm,headheight=14pt]{geometry}
\usepackage{fontspec}
\setmainfont{Latin Modern Roman}
\setsansfont{Latin Modern Sans}
\setmonofont{Latin Modern Mono}
\usepackage{amsmath,amssymb,mathtools}
\usepackage{unicode-math}
\setmathfont{Latin Modern Math}
\usepackage{microtype}
\usepackage{enumitem,xurl,needspace}
\usepackage[dvipsnames]{xcolor}
\usepackage{hyperref}
\hypersetup{unicode=true,colorlinks=true,linkcolor=MidnightBlue,urlcolor=MidnightBlue,
 pdftitle={500 Mathematical Problem Statements},pdfauthor={Source-annotated compilation},bookmarksnumbered=true}
\usepackage{bookmark}
\usepackage{tocloft}
\setlength{\cftsecnumwidth}{3em}
\cftsetpnumwidth{2.5em}
\cftsetrmarg{4em}
\usepackage{titlesec}
\titleformat{\section}{\Large\bfseries\raggedright}{\thesection}{0.7em}{}
\usepackage{fancyhdr}
\pagestyle{fancy}\fancyhf{}
\fancyhead[L]{\small\sffamily Mathematical problem statements}
\fancyhead[R]{\small Source edition 22}
\fancyfoot[C]{\thepage}
\setlength{\parindent}{0pt}
\setlength{\parskip}{5pt plus 1pt}
\setlength{\emergencystretch}{3em}
\setcounter{tocdepth}{1}
\setcounter{secnumdepth}{1}
\setlist[itemize]{leftmargin=3em,itemsep=5pt,topsep=4pt,parsep=0pt}
\allowdisplaybreaks
\newcommand{\N}{\mathbb{N}}
\newcommand{\Z}{\mathbb{Z}}
\newcommand{\Q}{\mathbb{Q}}
\newcommand{\R}{\mathbb{R}}
\newcommand{\C}{\mathbb{C}}
\newcommand{\F}{\mathbb{F}}
\newcommand{\A}{\mathbb{A}}
\newcommand{\PP}{\mathbb P}
\newcommand{\E}{\mathbb{E}}
\newcommand{\eps}{\varepsilon}
\newcommand{\dd}{\,\mathrm d}
\newcommand{\sref}[2]{\hyperlink{source:#1:#2}{[S#2]}}
\newcommand{\eref}[2]{\hyperlink{extra:#1:#2}{[E#2]}}
% Turn the next switch off for a shorter reading copy. The quotations preserve
% every exactTarget field, including qualifications omitted from a short summary.
\newif\ifshowcatalogue
\showcataloguetrue
\newcommand{\cataloguescope}[1]{\ifshowcatalogue
 \paragraph{Exact scope in the supplied catalogue.}
 \begin{quote}\small #1\end{quote}\fi}

% Standalone research-notebook wrapper; the source preamble above is retained.
\numberwithin{equation}{section}
\hypersetup{pdftitle={Research attempt -- problem 30006757},
 pdfauthor={Research notebook; original compilation retained}}
\fancyhead[L]{\small\sffamily Research notebook}
\fancyhead[R]{\small\sffamily Problem 30006757 | 27 September 2026}
\begin{document}
\setcounter{section}{0}
% ================================================================
% problemId: 30006757
% Canonical problem ID: 30006757
% exactTarget SHA-256: 941cb888427dc36af46d7a9542b5fa58014b8f56e714923f119f3c34e4aba2de
\clearpage
\section*{\texorpdfstring{Density of Hyperbolicity for Rational Maps}{Density of Hyperbolicity for Rational Maps}}\label{30006757}
\subsection{Definitions and mathematical statement}
Represent a degree-$d$ rational map by $f(z)=P(z)/Q(z)$ with coprime polynomials, modulo common nonzero scalar multiplication, giving the parameter space $\operatorname{Rat}_d\subset\PP^{2d+1}(\C)$. A critical point is a point of the sphere where the local mapping degree exceeds one. A periodic cycle is attracting if the modulus of its return-map multiplier is less than one. Call $f$ hyperbolic when every critical orbit converges to an attracting cycle. For each $d\ge2$, ask whether
\[
 \overline{\{f\in\operatorname{Rat}_d:f\text{ is hyperbolic}\}}=\operatorname{Rat}_d
\]
in the usual complex parameter topology. The target is density, not that every map is hyperbolic or structurally stable.
\par\smallskip\textit{Formulation and concept references:} \sref{30006757}{1}.
\cataloguescope{Determine whether hyperbolic rational maps, whose critical points are all attracted to attracting periodic cycles, are dense in the space of rational maps of the Riemann sphere of each fixed degree d >= 2.}
\subsection{Short English statement}
Can every rational dynamical system of a fixed degree be approximated arbitrarily closely by one whose critical points all settle into attracting cycles?
% BEGIN RESEARCH ADDITION ATTEMPT-155
\subsection{Research attempt}

\paragraph{Current frontier.}
Structural stability and hyperbolicity must not be identified: the density
of stable parameters does not by itself exclude stable non-hyperbolic
regions. Benini discusses rigidity, invariant line fields, and the special
role of the flexible Latt\`es family \sref{30006757}{1}, Sections 3 and 6.
A strong result directly relevant to simultaneous critical capture is
Aspenberg's Theorem A: a rational Misiurewicz map which is not flexible
Latt\`es can be approximated by hyperbolic maps; when its Julia set is the
whole sphere, all critical points of the approximants can be made periodic
\eref{30006757}{1}. Neither hypothesis is imposed in the target. The searches
for 2024--2026 work used for this entry did not identify a verified theorem
covering all of $\operatorname{Rat}_d$.

\paragraph{Attack route.}
An invariant-line-field route needs a rigidity theorem in precisely the
rational-map setting, including the exceptional families; a
renormalization route needs the necessary a priori bounds for every
recurrent configuration. Here we pursue a more quantitative route: make the
remaining critical points periodic by solving simultaneous critical-return
equations. The missing lemma is not simply transversality. One needs
transversality \emph{with control of the second derivatives on a sufficiently
large parameter neighborhood}. We prove a criterion, test it exactly in a
small instance, and then show that the most tempting derivative-growth
argument fails even at a familiar parabolic parameter.

\paragraph{Attempt: a simultaneous closing criterion.}
Work in a holomorphic parameter slice $\lambda\mapsto f_\lambda$ inside
$\operatorname{Rat}_d$, with $\lambda\in\C^m$ and center $0$. Assume that
all critical points are simple and labeled holomorphically on the
neighborhood under consideration. Designate $m$ critical points
$c_1(\lambda),\ldots,c_m(\lambda)$ as the ones not yet captured. For each
other critical point, assume that at $0$ a finite part of its orbit enters
a strictly trapping neighborhood of an attracting cycle. After shrinking
the parameter neighborhood, the same is true throughout it. This
persistence follows from the strict inclusion defining the trap and
continuity of the finitely many initial iterates.

Choose integers $n_i\ge1$ and compatible local coordinates near the
initial and return points, and set
\[
 R_i(\lambda)=f_\lambda^{n_i}(c_i(\lambda))-c_i(\lambda),
 \qquad R=(R_1,\ldots,R_m).
\]
The displayed subtraction is in those coordinates; their domains must
contain all points used on the closed parameter ball. Suppose
$A=DR(0)$ is invertible. For a norm on $\C^m$ and its induced operator
norms, write
\[
 \beta=\|A^{-1}R(0)\|,\qquad
 L=\sup_{\|\lambda\|\le2\beta}\|A^{-1}D^2R(\lambda)\|.
\]
\textbf{Closing lemma.} If $\beta>0$, the closed ball of radius $2\beta$
lies in the specified coordinate and trapping neighborhood, and
\begin{equation}\label{eq155-close}
 2L\beta\le\tfrac12,
\end{equation}
then some $\lambda_*$ with $\|\lambda_*\|\le2\beta$ makes $f_{\lambda_*}$
hyperbolic. If $\beta=0$, the same conclusion holds at $0$ under the other
assumptions.

\textit{Proof.} For $T(\lambda)=\lambda-A^{-1}R(\lambda)$, the mean value
formula along a segment gives
\[
 \|DT(\lambda)\|
 =\|A^{-1}(DR(0)-DR(\lambda))\|
 \le L\|\lambda\|\le\tfrac12.
\]
Thus $T$ is a contraction on that ball. Moreover
$\|T(0)\|=\beta$, so
$\|T(\lambda)\|\le\beta+\|\lambda\|/2\le2\beta$.
The contraction theorem supplies a fixed point, and invertibility of $A$
gives $R(\lambda_*)=0$. Each marked critical point is periodic, with period
dividing $n_i$, so its cycle is superattracting. All other critical points
remain in the preserved attracting basins. Hence every critical orbit is
attracted to an attracting cycle, exactly as required. \hfill$\square$

For example, in a Euclidean norm, a sufficient form of
\eqref{eq155-close} is
\[
 4H\|R(0)\|\le s^2,
 \qquad s=\sigma_{\min}(DR(0)),\qquad
 H\ge\sup_{\|\lambda\|\le2\beta}\|D^2R(\lambda)\|.
\]
Both the size of the coordinate neighborhood and the trapping radii remain
separate conditions. The dimension of the slice, independence of its
critical-return equations, and a lower bound on $s$ are not automatic.

\paragraph{Attempt: an exact arithmetic test.}
For $f_c(z)=z^2+c$, infinity is already a superattracting fixed point, so
only the finite critical point $0$ must be closed. The third-return
polynomial is
\[
 R_3(c)=f_c^3(0)=c^4+2c^3+c^2+c.
\]
Take
\[
 c_0=-\frac{17549}{10000},\qquad r=\frac1{10000}.
\]
Exact rational calculation gives
\[
 R_3(c_0)=\frac{1261775714801}{10^{16}},\qquad
 R_3'(c_0)=-\frac{1412458905149}{250000000000}.
\]
On $|c-c_0|\le r$, one has $|c|<2$, and hence
$|R_3''(c)|=|12c^2+12c+2|\le74$. With $A=R_3'(c_0)$,
\[
 \beta=|R_3(c_0)/A|<r/2,\qquad q=74r/|A|<1/2.
\]
The same contraction proof, now on the fixed radius-$r$ disk, produces a
zero of $R_3$ in that disk. Since the disk excludes $c=0$, the critical
point has exact period three: the only proper positive divisor of three is
one, and $f_c(0)=0$ would require $c=0$. This is a finite certificate of a
hyperbolic parameter, not a density argument. The attached checking script
verifies the rational inequalities without floating-point decisions.

\paragraph{Attempt: a parabolic stress calculation.}
A natural proposed general lemma would say that a sufficiently large
parameter derivative makes the critical-return equation solvable nearby.
The following calculation shows why that is not enough for the preceding
criterion. Put $z_0(c)=0$ and $z_{n+1}(c)=z_n(c)^2+c$. At $c=1/4$, let
\[
 x_n=\tfrac12-z_n(1/4),\quad
 a_n=z_n'(1/4),\quad b_n=z_n''(1/4).
\]
Direct differentiation, with no interchange of infinite limits, gives
\begin{align}
 x_{n+1}&=x_n-x_n^2, & x_0&=\tfrac12,\label{eq155-x}\\
 a_{n+1}&=(1-2x_n)a_n+1, & a_0&=0,\label{eq155-a}\\
 b_{n+1}&=(1-2x_n)b_n+2a_n^2, & b_0&=0.\label{eq155-b}
\end{align}
The sequence $x_n$ decreases to zero, since a nonnegative limit in
\eqref{eq155-x} must have square zero. Moreover
\[
 \frac1{x_{n+1}}-\frac1{x_n}=\frac1{1-x_n}\longrightarrow1,
\]
so the Ces\`aro theorem applied to the telescoping sum yields $nx_n\to1$.

For completeness, use the elementary recurrence fact
\begin{equation}\label{eq155-recurrence}
 y_{n+1}=\left(1-\frac{\alpha}{n}+o(n^{-1})\right)y_n
       +c n^k(1+o(1))
 \quad\Longrightarrow\quad
 y_n\sim\frac{c}{\alpha+k+1}n^{k+1},
\end{equation}
when $\alpha,c>0$, $k\ge0$, and $y_n\ge0$ eventually. Here is a proof
rather than an appeal to a differential-equation approximation. Substitute
$C n^{k+1}$ in the recurrence: its increment minus the right-hand increment
is $((\alpha+k+1)C-c)n^k+o(n^k)$. Constants just above and below
$c/(\alpha+k+1)$ therefore give eventual super- and subsolutions. An
initial discrepancy is multiplied by products of eventually positive
coefficients at most $1-\alpha/(2n)$; those products tend to zero, and their
contribution divided by $n^{k+1}$ tends to zero. Comparison and then letting
the two constants approach their common limit proves
\eqref{eq155-recurrence}.

Equations \eqref{eq155-a} and \eqref{eq155-b} now imply
\[
 a_n\sim\frac n3,\qquad b_n\sim\frac{2n^3}{45}.
\]
For the equation $R_n(c)=z_n(c)=0$, at the parabolic center this gives
\begin{equation}\label{eq155-failure}
 \left|\frac{R_n(1/4)}{R_n'(1/4)}\right|\sim\frac{3}{2n}
 \longrightarrow0,
 \qquad
 \frac{|R_n(1/4)R_n''(1/4)|}{|R_n'(1/4)|^2}\sim\frac n5
 \longrightarrow\infty.
\end{equation}
Thus even the \emph{pointwise} lower bound on the normalized Hessian
already prevents \eqref{eq155-close} for large $n$. A small first Newton
step conceals large distortion. This is a proved failure of this sufficient
criterion at these centers, not a proof that the equations have no nearby
roots.

\paragraph{Stress test.}
The parabolic example is certainly not a counterexample to density. For
$0<c<1/4$, the smaller fixed point
$\alpha_c=(1-\sqrt{1-4c})/2$ has multiplier $2\alpha_c<1$.
Starting at $0$, iteration increases monotonically and stays below
$\alpha_c$, so it converges to that fixed point. Together with the
attracting point at infinity, this proves hyperbolicity for these maps,
which approach $c=1/4$. The failure in \eqref{eq155-failure} therefore
identifies a weakness of the chosen proof device, not of the target.

One must also avoid requiring \emph{all} critical points to be periodic in
every neighborhood, even inside an attracting component: attracting
orbits need not become periodic under arbitrary small perturbation. Our
criterion leaves already attracted points alone. Multiple critical points
can be avoided when first choosing a nearby generic parameter, but that
does not establish the uniform return or distortion estimates needed at
the next stage. Likewise, a closing result for one critical orbit does not
control all the other critical orbits. Aspenberg's simultaneous control
uses its nonrecurrence hypotheses \eref{30006757}{1}; they cannot be removed
by just applying the inverse function theorem repeatedly.

\paragraph{Outcome.}
\textbf{Outcome: Exploratory approach with unresolved gap.}

The attempt yields a fully proved quantitative simultaneous-closing lemma,
an exact rational period-three certificate, and a parabolic asymptotic
obstruction to the naive derivative-growth implementation. These are
lemmas and tests within the attempt, not a claim of a new density theorem
or of priority for contraction-based closing methods.

A complete proof along this route would have to show that every parameter
neighborhood contains a center and a collection of critical-return equations
satisfying the radius, transversality, and distortion conditions, or replace
this criterion with one that survives the demonstrated parabolic failure.
No such uniform statement is proved, especially for recurrent critical
orbits and the exceptional flexible families. The general density problem
is not resolved here.
% END RESEARCH ADDITION ATTEMPT-155
\subsection{Sources}
\begin{itemize}\raggedright
\item[\textnormal{[S1]}]\hypertarget{source:30006757:1}{} A. M. Benini, 'A survey on MLC, Rigidity and related topics', arXiv:1709.09869 (2017, revised 2018).
\url{https://arxiv.org/abs/1709.09869}.
{\small\textit{Catalogue formulation source.}}
% BEGIN RESEARCH ADDITION SOURCES-155
\item[\textnormal{[E1]}]\hypertarget{extra:30006757:1}{} Magnus Aspenberg,
\emph{Perturbations of rational Misiurewicz maps}, arXiv:0804.1106v2,
22 June 2009, Definition 1.1, Theorem A, and Sections 3--5.
\url{https://arxiv.org/html/0804.1106v2}.
{\small\textit{Consulted 27 September 2026. The exclusion of flexible
Latt\`es maps and the Misiurewicz hypothesis are essential to the cited result.}}
% END RESEARCH ADDITION SOURCES-155
\end{itemize}

\end{document}
